La realización de esta experiencia ha permitido cumplir satisfactoriamente los objetivos generales y específicos planteados, logrando una comprensión práctica y metódica de la operación manual del osciloscopio digital. Los resultados obtenidos confirman que la fiabilidad de cualquier medición en el dominio del tiempo depende de la interrelación de tres sistemas clave: el sistema de disparo (trigger), el acoplamiento de entrada y las escalas de visualización (ganancia y base de tiempo).

El hallazgo más significativo es la confirmación de que el sistema de trigger es el pilar fundamental para obtener mediciones fiables. Se comprobó experimentalmente que, para lograr una traza estable, el nivel de disparo debe estar configurado obligatoriamente dentro del rango de amplitud de la señal de referencia. Fuera de este rango, el modo Normal cesa la adquisición (estado "WAIT").

En el análisis de acoplamiento, se validó que el modo CC es esencial para medir la señal completa (AC+DC), obteniendo un valor medio de \SI{3.930}{\volt} para CH1. Se demostró que el modo CA filtra esta componente de offset sin alterar la relación de fase entre las señales.

La exploración de las escalas demostró que una ganancia (V/div) inadecuada conduce a errores graves: una sensibilidad muy baja vuelve la señal imperceptible (Figura 11), mientras que una sensibilidad muy alta provoca saturación, inutilizando la medición de amplitud (Figura 13) . El modo X-Y se validó como una técnica potente, no solo para visualizar la evolución del desfase (Figura 15), sino también para identificar relaciones de frecuencia armónicas (2:1, 1:2, 3:2) mediante el conteo de tangentes en las figuras de Lissajous (Figura 16).

Una observación de gran importancia fue la ineficacia de la sincronización con la red eléctrica (Line) para las señales generadas en el laboratorio. Se constató que la frecuencia de la red eléctrica fluctúa, mientras que la del generador de funciones es fija (\SI{50.00}{\hertz}). Esta mínima diferencia de frecuencias impide una sincronización estable, provocando un "desplazamiento" (drift) continuo de las trazas .

Como investigación futura, se desprende que el trigger de "Line" sería, por el contrario, la opción óptima en escenarios donde el circuito bajo análisis esté alimentado directamente por la red eléctrica. En ese caso, tanto la referencia del trigger como las señales a medir compartirían la misma frecuencia fundamental fluctuante, permitiendo una sincronización perfecta. Este principio será clave en la próxima experiencia de laboratorio (Experiencia 5), donde se analizan circuitos conectados a la red.

Finalmente, la captura de la señal transitoria (Figuras 17-19) demostró la indispensabilidad del modo de disparo Single y la importancia de una correcta configuración de la base de tiempo y la posición horizontal del trigger para aislar fenómenos específicos, como se evidenció en el intento fallido de aislar el flanco de bajada.